Un fil de cuivre est formé d'une section épaisse suivie d'une section mince, en série. Un courant constant le traverse. Dans la section mince, comment se comparent le courant et la vitesse de dérive des électrons à ceux de la section épaisse ?

\choices

\noindent\textbf{Justification :} \hrulefill \\[2.5ex]
\rule{\linewidth}{0.4pt} \\[2.5ex]
\rule{\linewidth}{0.4pt}
\vfill
